\documentclass[10pt]{article}

\usepackage[margin=0.65in]{geometry}
\usepackage[T1]{fontenc}
\usepackage[utf8]{inputenc}
\usepackage{lmodern}
\usepackage{hyperref}
\usepackage{longtable}
\usepackage{array}
\usepackage{enumitem}
\usepackage{parskip}
\usepackage{xcolor}
\usepackage{titlesec}

\definecolor{cpastblue}{HTML}{1D4ED8}
\definecolor{cpastgreen}{HTML}{047857}
\definecolor{cpastgold}{HTML}{B45309}
\definecolor{cpastink}{HTML}{111827}
\definecolor{softblue}{HTML}{EFF6FF}

\hypersetup{
  colorlinks=true,
  urlcolor=cpastblue,
  linkcolor=cpastblue
}

\titleformat{\section}{\large\bfseries\color{cpastink}}{\thesection}{0.6em}{}
\titleformat{\subsection}{\normalsize\bfseries\color{cpastink}}{\thesubsection}{0.6em}{}

\newcommand{\doclink}[2]{\href{run:#1}{#2}}

\begin{document}

\begin{center}
{\LARGE\bfseries CPAST Self-Evaluation}\\[0.25em]
{\large Piter Garcia --- Pine Brook Elementary IGNITE --- Spring 2026}\\[0.4em]
\colorbox{softblue}{\parbox{0.92\linewidth}{\centering
\textbf{Purpose:} Candidate self-evaluation draft for the CPAST three-way conference, grounded in local transcripts, lesson packets, and placement documents.
}}
\end{center}

\section{Quick Scan}

\begin{itemize}[leftmargin=*]
  \item \textbf{Rubric:} \doclink{../warner-school-cpast-rubric.docx}{Warner School CPAST rubric}
  \item \textbf{Evidence examples:} \doclink{../warner-cpast-examples-of-evidence.docx}{Warner CPAST examples of evidence}
  \item \textbf{Meeting prep:} \doclink{./cpast-three-way-conference-prep.pdf}{CPAST three-way conference prep packet}
  \item \textbf{Lesson hub:} \doclink{../../pine-brook-elementary/lessons/lessons-notes.md}{Pine Brook lesson hub}
  \item \textbf{Transcript hub:} \doclink{../../transcripts/transcripts-notes.md}{Teaching transcript hub}
\end{itemize}

\section{Overall Self-Assessment}

I would self-assess my overall performance as \textbf{meeting expectations, with several areas approaching exceeds expectations}, especially in digital tools, curriculum/resource design, collaboration, responsiveness to feedback, and inclusive planning. My strongest evidence comes from the way I helped build and revise computer science lessons across Minecraft Education, Sphero, Code.org, Book Creator, Tinkercad, Google Slides, Canva, and AI-supported project work.

The main growth area I would name honestly is \textbf{assessment documentation}. I used formative checks, observation, student work, and class adjustments, but I still need a cleaner routine for recording data, preserving de-identified examples, and connecting evidence directly to assessment decisions.

\section{Suggested CPAST Rating Table}

\small
\begin{longtable}{>{\raggedright\arraybackslash}p{0.35\linewidth} >{\centering\arraybackslash}p{0.12\linewidth} >{\raggedright\arraybackslash}p{0.45\linewidth}}
\textbf{CPAST Area} & \textbf{Rating} & \textbf{Evidence Summary} \\
\hline
A. Focus for Learning: Standards and Objectives/Targets & 3 & Lesson packets connect workbook expectations, CS standards, and student-facing targets; some lessons need cleaner final alignment. \\
B. Materials and Resources & 4 & Strong use of varied materials across digital tools, robotics, Minecraft, AI, slides, rubrics, and lesson packets. \\
C. Assessment of P-12 Learning & 3 & Rubrics, exit tickets, observation, and artifacts are present; documentation should become more systematic. \\
D. Differentiated Methods & 3 & Visuals, chunking, modeling, partner structures, and audience profiles support access; individual evidence can be stronger. \\
E. Learning Target and Directions & 3 & Student-facing directions and slides improved over time; early materials sometimes needed clearer directions. \\
F. Critical Thinking & 3 & Students engaged in debugging, design choices, AI evaluation, and problem solving; explicit reflection prompts can improve this. \\
G. Checking for Understanding and Adjusting Instruction & 3 & Transcripts show adjustment, troubleshooting, and re-explanation; formal data capture needs strengthening. \\
H. Digital Tools and Resources & 4 & Clear strength: technology was central, developmentally appropriate, and tied to CS learning. \\
I. Safe and Respectful Learning Environment & 3 & Evidence shows routines, safety, digital citizenship, and student support; next step is stronger proactive whole-room systems. \\
J. Data-Guided Instruction & 2--3 & Informal data and observations guided decisions, but a stronger record-keeping system is needed. \\
K. Feedback to Learners & 3 & Feedback occurred during projects, coding, and revision; written or captured feedback should be more consistent. \\
L. Assessment Techniques & 3 & Rubrics, exit tickets, project artifacts, and performance tasks are present; diagnostic documentation can be stronger. \\
M. Connections to Research and Theory & 3--4 & Strong evidence through UDL, CS standards, AI literacy, digital citizenship, and reflective improvement. \\
N. Participates in Professional Development & 3 & Coursework, professional conversations, artifact protocols, and feedback cycles informed practice. \\
O. Communication with Parents/Guardians & N/A or 3 & Parent communication was routed through Derek according to placement policy. \\
P. Punctuality & 3 & No contrary evidence found; schedule commitments were documented and followed as part of the placement plan. \\
Q. Meets Deadlines and Obligations & 3 & Organized documentation, lesson materials, and required paperwork are present; earlier organization is a growth point. \\
R. Preparation & 3 & Lesson packets, slides, rubrics, references, and materials show preparation; some packets needed later cleanup. \\
S. Collaboration & 4 & Strong evidence of co-planning, feedback conversations, and coordination with Derek, Zenon, Sue, and school context. \\
T. Advocacy for Learners or Profession & 3--4 & Evidence includes inclusion, AI safety, access, scaffolding, and student independence. \\
U. Responds Positively to Feedback & 4 & Repeated revision based on feedback, especially around rubrics, AI lessons, inclusion, and lesson clarity. \\
\end{longtable}
\normalsize

\section{Evidence-Based Reflection By Domain}

\subsection{Planning for Instruction and Assessment}

\textbf{Self-evaluation:} I meet expectations in planning and am moving toward exceeds in materials/resources. I built and organized lesson packets across grades, tied them to the Pine Brook IGNITE curriculum, added standards and references where available, and created student-facing lesson materials.

\textbf{Evidence to cite:}
\begin{itemize}[leftmargin=*]
  \item \doclink{../../pine-brook-elementary/teaching-placement-ignite-curriculum-gcsd-25-26/teaching-placement-ignite-curriculum-gcsd-25-26-notes.md}{IGNITE curriculum notes}
  \item \doclink{../../pine-brook-elementary/lessons/lesson-outline-notes.md}{Lesson outline map}
  \item \doclink{../../pine-brook-elementary/lessons/grade-4-baris-fowler-schrank-autore/minecraft-coding-fundamentals/minecraft-coding-fundamentals.md}{Minecraft coding fundamentals}
  \item \doclink{../../pine-brook-elementary/lessons/grade-3-lamanaco-lallucci-regelsberger-tandoi/magicschool-student-choice-ai-partner-project/magicschool-student-choice-ai-partner-project.md}{MagicSchool student choice AI partner project}
\end{itemize}

\textbf{Growth note:} Planning became stronger when I separated official sources, workbook expectations, transcript evidence, and enriched lesson plans. I still need that structure to be more consistent for every lesson.

\subsection{Instructional Delivery}

\textbf{Self-evaluation:} I meet expectations in instructional delivery, with a clear strength in using technology to support CS learning. The transcripts show teaching around coding, robotics, AI, Minecraft Education, Sphero, Tinkercad, Book Creator, and Google tools, along with real-time troubleshooting and adjustment.

\textbf{Evidence to cite:}
\begin{itemize}[leftmargin=*]
  \item \doclink{../../transcripts/grade-4/260313-grade-4-baris-minecraft-fundamentals.md}{Grade 4 Minecraft fundamentals transcript}
  \item \doclink{../../transcripts/grade-4/260302-grade-4-baris-minecraft-test.md}{Grade 4 Minecraft test lesson transcript}
  \item \doclink{../../transcripts/grade-3/260227-grade-3-lamanaco-sphero-loops.md}{Grade 3 Sphero loops transcript}
  \item \doclink{../../transcripts/grade-5/260304-grade-5-pickett-tinkercad.md}{Grade 5 Tinkercad transcript}
  \item \doclink{../../transcripts/grade-2/260211-grade-2-callon-ai-oceans.md}{Grade 2 AI for Oceans transcript}
\end{itemize}

\textbf{Growth note:} I need to keep student-facing directions shorter and more visual. Slide revisions show that I recognized when directions were too teacher-facing or overloaded and began correcting them into clearer student language.

\subsection{Assessment}

\textbf{Self-evaluation:} I meet expectations in assessment design, but assessment documentation is my clearest next growth step. I used rubrics, exit tickets, project products, observation, and performance tasks, especially in AI and project-based lessons. I need a more consistent routine for recording formative data and linking it directly to instructional decisions.

\textbf{Evidence to cite:}
\begin{itemize}[leftmargin=*]
  \item \doclink{../../pine-brook-elementary/lessons/grade-3-lamanaco-lallucci-regelsberger-tandoi/magicschool-student-choice-ai-partner-project/lesson1/lesson1-exit_ticket.md}{Lesson 1 AI exit ticket}
  \item \doclink{../../pine-brook-elementary/lessons/grade-3-lamanaco-lallucci-regelsberger-tandoi/magicschool-student-choice-ai-partner-project/lesson1/lesson1-rubric.md}{Lesson 1 AI rubric}
  \item \doclink{../../pine-brook-elementary/lessons/grade-3-lamanaco-lallucci-regelsberger-tandoi/magicschool-student-choice-ai-partner-project/lesson2/lesson2-rubric.md}{Lesson 2 AI rubric}
  \item \doclink{../warner-cpast-examples-of-evidence.docx}{CPAST examples of evidence}
\end{itemize}

\textbf{Growth note:} Bring one or two concrete assessment examples to the meeting and name this as a real improvement goal: more systematic documentation of checks for understanding, student work samples, and instructional adjustments.

\subsection{Analysis of Teaching}

\textbf{Self-evaluation:} I meet expectations and may approach exceeds because I repeatedly connected lesson design to theory, standards, and feedback. I used UDL, CS standards, AI literacy guidance, digital citizenship, official lesson sources, and feedback to revise materials.

\textbf{Evidence to cite:}
\begin{itemize}[leftmargin=*]
  \item \doclink{../../transcripts/non-lesson-context/260210-placement-expectations-and-observation-planning.md}{Placement expectations and observation planning}
  \item \doclink{../../pine-brook-elementary/lessons/grade-3-lamanaco-lallucci-regelsberger-tandoi/magicschool-student-choice-ai-partner-project/magicschool-student-choice-ai-partner-project.bib}{MagicSchool project bibliography}
  \item \doclink{../../transcripts/non-lesson-context/260317-caregiving-sel-and-family-support-context.md}{Caregiving SEL and family support context}
  \item \doclink{../../transcripts/non-lesson-context/260311-student-dependency-and-return-planning.md}{Student dependency and return planning}
\end{itemize}

\textbf{Growth note:} The next step is to make reflection more routine and concise: after each major lesson, document what happened, what evidence showed, what changed, and what should change next.

\subsection{Professional Commitment, Relationships, and Feedback}

\textbf{Self-evaluation:} I meet expectations in professional commitment and would rate collaboration and responsiveness to feedback as strengths. The evidence shows collaboration with Derek around lesson selection, classroom routines, observation opportunities, inclusion, and instructional decisions. It also shows responsiveness to feedback around rubrics, AI lessons, and lesson clarity.

\textbf{Evidence to cite:}
\begin{itemize}[leftmargin=*]
  \item \doclink{../expectations/letter-of-expectations-spring-2026-work.md}{Letter of Expectations work file}
  \item \doclink{./student-teaching-internship-action-items-and-key-information.md}{Student teaching internship action items}
  \item \doclink{../../transcripts/non-lesson-context/260210-placement-expectations-and-observation-planning.md}{Placement expectations and observation planning}
  \item \doclink{../../transcripts/non-lesson-context/260311-student-dependency-and-return-planning.md}{Student dependency and return planning}
\end{itemize}

\textbf{Growth note:} My next step is to make collaboration artifacts easier to show: short notes from planning conversations, feedback received, and how I changed the lesson afterward.

\section{Three Talking Points For The Meeting}

\begin{enumerate}[leftmargin=*]
  \item \textbf{My strongest growth was turning technology-rich activities into clearer CS instruction.} Early lessons often involved tool access and troubleshooting; later materials show stronger links between tools, standards, student-facing directions, and evidence of learning.
  \item \textbf{Inclusion became a planning lens, not an add-on.} The class audience profiles, transcript notes, and lesson scaffolds show increasing attention to visual supports, chunked steps, partner structures, emotional regulation, and gradual release.
  \item \textbf{My next professional goal is cleaner assessment documentation.} I can point to rubrics, exit tickets, project artifacts, and observed adjustments, but I need a more consistent system for collecting and explaining formative data.
\end{enumerate}

\section{Suggested Growth Goals}

\begin{longtable}{>{\raggedright\arraybackslash}p{0.26\linewidth} >{\raggedright\arraybackslash}p{0.34\linewidth} >{\raggedright\arraybackslash}p{0.32\linewidth}}
\textbf{Growth Goal} & \textbf{Why It Matters} & \textbf{Concrete Next Step} \\
\hline
Build a simple formative assessment log & Makes data-guided instruction easier to prove & After each lesson, record one check, what it showed, and what changed. \\
Keep lesson plans in the three-layer model & Prevents lesson folders from becoming generic summaries & Separate official lesson, provided/base plan, and enriched transcript-informed plan. \\
Preserve de-identified student evidence & Makes assessment and feedback concrete & Save one artifact, screenshot, exit ticket, or observation note per major lesson. \\
Tighten student-facing directions & Supports access and reduces confusion & Use fewer words, visual checkpoints, and ``I can'' success criteria. \\
\end{longtable}

\section{Meeting-Friendly Closing Statement}

My biggest takeaway is that effective CS teaching is not just about giving students access to interesting tools. The work is in making the learning target visible, designing supports so students can enter the task, checking whether they understand, and then adjusting when the evidence says they need something different. I have grown in building those materials and responding to feedback, and my next step is to make assessment evidence more systematic and easier to show.

\end{document}
